\section{Compilation Pipeline}
\label{sec:pipeline}

The compiler is a straight pipeline from source to artifacts:

\begin{center}
\begin{tikzpicture}[
  every node/.style={font=\small},
  box/.style={draw, rounded corners, inner sep=4pt, minimum width=2.4cm,
              minimum height=6mm, fill=blue!6, align=center},
  ir/.style ={draw, rounded corners, inner sep=4pt, minimum width=2.4cm,
              minimum height=6mm, fill=yellow!20, align=center},
  arr/.style={-{Stealth[length=2mm]}, thick}
]
\node[box] (src) at (0,0) {Source\\\TT{.vg}};
\node[box] (ast) at (2.6,0) {AST};
\node[box] (mac) at (5.2,0) {Macro-\\inlined AST};
\node[box] (chk) at (7.8,0) {Type-\\checked};
\node[ir]  (eg)  at (10.4,0) {EventGraph};

\node[box] (sol) at (4.2,-1.6) {Solidity / Vyper};
\node[box] (gam) at (7.0,-1.6) {Gambit EFG};
\node[box] (smt) at (9.5,-1.6) {SMT-LIB};
\node[box] (scr) at (11.8,-1.6) {Scribble};
\node[box] (gal) at (4.2,-2.4) {Gallina / Lean};
\node[box] (mai) at (7.0,-2.4) {MAID JSON};
\node[box] (btc) at (9.5,-2.4) {Bitcoin / LN};
\node[box] (gv)  at (11.8,-2.4) {GraphViz};

\draw[arr] (src) -- (ast);
\draw[arr] (ast) -- (mac);
\draw[arr] (mac) -- (chk);
\draw[arr] (chk) -- (eg);

\foreach \t in {sol, gam, smt, scr, gal, mai, btc, gv}
  \draw[arr, gray] (eg.south) to[bend left=8] (\t.north);
\end{tikzpicture}
\end{center}

The frontend (\file{vegas/frontend/}) parses with an ANTLR-generated
parser, inlines macros, and runs a four-pass type checker that
validates type aliases, macro signatures, the protocol's
\TT{join}/\TT{yield}/\TT{reveal}/\TT{commit}/\TT{withdraw} structure,
and the visibility discipline of guards.  The frontend then lowers
to an intermediate representation \TT{GameIR}, which packages a
finite EventGraph together with a per-role payoff expression and
the chance-role set.  The EventGraph is the spine: every backend
reads it directly and emits its own artifact without revisiting the
surface syntax.

Two transformation passes operate on the IR.  The first builds the
EventGraph from typechecked AST: it groups simultaneous moves into
phases, infers dependencies (data, commit-reveal, phase-order), and
computes per-node visibility (\COMMIT, \REVEAL, or \PUBLIC).  The
second is the \emph{commit-reveal expansion} pass
(\S\ref{sec:commit-reveal}), which detects concurrent
public-write nodes (the ``risk partners'') and rewrites each such
node into a commit-then-reveal pair, with predecessors arranged so
that no reveal is enabled while a sibling's commit is still
unfinished.

After the rewrite, the EventGraph satisfies three invariants:
acyclicity; commit-before-reveal (for each pair of \COMMIT/\REVEAL
nodes on the same field, the commit dominates the reveal in the
DAG); and read-visibility (every field referenced by a guard is
either public at the action's point in the DAG, or has already been
revealed in a predecessor).  The graph constructor refuses to
return an EventGraph that violates these invariants; backends can
assume them.

Backends fall into two groups.  The \emph{operational} backends ---
Solidity, Vyper, Bitcoin/Lightning --- emit an executable artifact
that implements the protocol on its target runtime.  The
\emph{analytical} backends --- Gambit, SMT-LIB, MAID, Scribble,
Gallina/Lean --- emit a non-executable representation suitable for
external tools (Gambit for equilibrium computation, Z3 for
reachability, Coq/Lean for proof, etc.).  Both groups consume the
same EventGraph.  The next several sections walk through each
in turn.
